% Part:sets-functions-relations
% Chapter: size-of-sets
% Section: comparing-sizes

\documentclass[../../../include/open-logic-section]{subfiles}

\begin{document}

\olfileid{sfr}{siz}{car}

\olsection{वेगवेगळ्या आकारमानांचे संच आणि कँटर यांचे प्रमेय}

\begin{explain}
दोन संचांचे आकारमान समान असते या कल्पनेचे नेमके विधान आपण दिले आहे.
एक संच दुसऱ्यापेक्षा लहान असतो या कल्पनेचेही नेमके विधान करता येते.
``लहान (किंवा तुल्यबल) असणे'' याच्या व्याख्येत संचांदरम्यानचे
!!a{bijection} न घेता पहिल्या संचापासून दुसऱ्या संचाकडे जाणारे
!!a{injection} घ्यावे लागेल. असे फलन अस्तित्वात असेल, तर पहिल्या
संचाचे आकारमान दुसऱ्या संचाच्या आकारमानापेक्षा लहान किंवा त्याच्याइतके
असते. अंतःप्रज्ञेने, एका संचापासून दुसऱ्या संचाकडे जाणारे
!!a{injection} हे फलनाच्या व्याप्तीत प्रांताइतके तरी !!{element}s आहेत
याची खात्री देते, कारण प्रांतातील कोणतेही दोन !!{element}s व्याप्तीतील
एकाच !!{element} कडे पाठवले जात नाहीत.
\end{explain}

\begin{defn}
$A$ हा $B$ पेक्षा \emph{आकारमानाने मोठा नाही}, हे $\cardle{A}{B}$ असे
लिहितो, तेव्हा आणि केवळ तेव्हाच, जेव्हा $f \colon A \to B$ असे
!!a{injection} असते.
\end{defn}

हा संबंध परावर्ती आणि संक्रमक आहे, पण सममित नाही, हे स्पष्ट आहे (हे
सरावासाठी सोडले आहे). एक संच दुसऱ्या संचापेक्षा (काटेकोरपणे) लहान आहे
असे सांगणारी कल्पनाही मांडता येते.

\begin{defn}
$A$ हा $B$ पेक्षा \emph{आकारमानाने लहान} आहे, हे $\cardless{A}{B}$ असे
लिहितो, तेव्हा आणि केवळ तेव्हाच, जेव्हा $f\colon A \to B$ असे
!!a{injection} असते, पण $g\colon A
\to B$ असे कोणतेही !!{bijection} नसते; म्हणजे, $\cardle{A}{B}$ आणि
$\cardneq{A}{B}$.
\end{defn}

हा संबंध अपरावर्ती आणि संक्रमक आहे, हे स्पष्ट आहे. (हे सरावासाठी
सोडले आहे.) या संकेतनाने संच $A$ हा !!{enumerable} असतो तेव्हा आणि
केवळ तेव्हाच $\cardle{A}{\Nat}$, आणि $A$ हा !!{nonenumerable} असतो
तेव्हा आणि केवळ तेव्हाच $\cardless{\Nat}{A}$, असे म्हणता येते. त्यामुळे
\oliflabeldef{sfr:siz:nen-alt:thm:nonenum-pownat}{%
\olref[sfr][siz][nen-alt]{thm:nonenum-pownat}
हे प्रमेय $\cardless{\Nat}{\Pow{\Nat}}$ असे पुन्हा मांडता येते}{%
\olref[sfr][siz][nen]{thm:nonenum-pownat}
हे प्रमेय $\cardless{\PosInt}{\Pow{\PosInt}}$ असे पुन्हा मांडता येते}.
खरे तर, हा शेवटचा मुद्दा \emph{पूर्णपणे व्यापक} आहे, हे
\citet{Cantor1892} यांनी सिद्ध केले:

\begin{thm}[कँटर]\ollabel{thm:cantor}
कोणत्याही संच $A$ साठी $\cardless{A}{\Pow{A}}$.
\end{thm}

\begin{proof}
$f(x) = \{x\}$ हे $f \colon A \to \Pow{A}$ असे !!a{injection} आहे;
कारण $x \neq y$ असेल, तर विस्तारात्मकतेमुळे $\{x\} \neq \{y\}$ सुद्धा
असते, आणि म्हणून $f(x) \neq f(y)$. त्यामुळे $\cardle{A}{\Pow{A}}$.

\begin{editorial}
\olref[nen]{sec} समाविष्ट असेल, तर आपण संथ सिद्धता देतो; अन्यथा
\olref[nen-alt]{sec} शी जुळणारी अधिक जलद सिद्धता देतो.
\end{editorial}

\oliflabeldef{sfr:siz:nen:sec}{% 
  आता $g\colon A \to
  \Pow{A}$ असे !!a{surjective} फलन, आणि त्यामुळे !!a{bijective} फलनही,
  असू शकत नाही हे दाखवू; म्हणून $\cardneq{A}{\Pow{A}}$. त्यासाठी
  $g\colon A \to \Pow{A}$ असे समजा. $g$ हे सर्वत्र परिभाषित असल्याने
  प्रत्येक $x \in A$ हा $g(x)
  \subseteq A$ अशा उपसंचाकडे पाठवला जातो. $g$ हे आच्छादक असू शकत नाही,
  हे दाखवता येते. त्यासाठी व्याख्येमुळेच~$g$ च्या व्याप्तीत येऊ न
  शकणारा $\overline{A} \subseteq A$ असा उपसंच परिभाषित करू. पुढील संच घ्या:
  \[
  \overline{A} = \Setabs{x \in A}{x \notin g(x)}.
  \]
  प्रत्येक $x \in A$ साठी $g(x)$ परिभाषित असल्याने $\overline{A}$ हा
  स्पष्टपणे~$A$ चा सुयोग्य रीतीने परिभाषित उपसंच आहे. पण तो~$g$ च्या
  व्याप्तीत असू शकत नाही. $x \in A$ स्वैर असू द्या; आपण
  $\overline{A} \neq g(x)$ दाखवू. $x \in g(x)$ असेल, तर तो
  $x \notin g(x)$ हा गुणधर्म पूर्ण करत नाही, आणि म्हणून~$\overline{A}$
  च्या व्याख्येवरून $x \notin
  \overline{A}$. $x \in \overline{A}$ असेल, तर त्याने~$\overline{A}$
  चा व्याख्यात्मक गुणधर्म पूर्ण केला पाहिजे; म्हणजे $x \in A$ आणि
  $x \notin
  g(x)$. $x$ स्वैर असल्याने प्रत्येक $x \in A$ साठी $x \in g(x)$ तेव्हा
  आणि केवळ तेव्हाच $x \notin \overline{A}$; म्हणून
  $g(x) \neq \overline{A}$.
  दुसऱ्या शब्दांत, $\overline{A}$ हा~$g$ च्या व्याप्तीत येऊ शकत नाही;
  त्यामुळे~$g$ आच्छादक आहे या गृहीतकाला व्याघात होतो.}{आता
  $\cardneq{A}{\Pow{A}}$ दाखवणे बाकी आहे. व्याघातासाठी
  $\cardeq{A}{\Pow{A}}$, म्हणजे $g \colon A \to \Pow{A}$ असे काही
  !!{bijection} आहे, असे समजा. आता पुढील संच विचारात घ्या:
  \[
    D = \Setabs{x \in A}{x \notin g(x)}
  \]
  $D \subseteq A$ आहे, म्हणून $D \in \Pow{A}$. $g$ हे
  !!a{bijection} असल्यामुळे $g(y) = D$ असा एखादा $y \in A$ घटक आहे.
  पण आता:
  \[
    y \in g(y) \text{ तेव्हा आणि केवळ तेव्हाच } y \in D \text{ तेव्हा आणि केवळ तेव्हाच } y \notin g(y).
  \]
  हा व्याघात आहे; म्हणून $\cardneq{A}{\Pow{A}}$.}{}
%READERNOTE{OLSIZ-007: गोठवलेल्या संथ सिद्धतेत x स्वैरपणे A मधून घेतल्यानंतर निष्कर्षाचा आवाका चुकून “प्रत्येक x जो A-bar मध्ये आहे” असा दिला आहे. g(x) आणि A-bar वेगळे आहेत हे प्रत्येक x∈A साठी दाखवणे आवश्यक आहे; मराठी सिद्धतेत तो आवाका दुरुस्त केला आहे.}
\end{proof}

\begin{explain}
\oliflabeldef{sfr:siz:nen:thm:nonenum-pownat}{
  \olref{thm:cantor} ची सिद्धता आणि
  \olref[nen]{thm:nonenum-pownat} ची सिद्धता यांची तुलना करणे उद्बोधक
  आहे. तेथे~$\PosInt$ च्या उपसंचांची कोणतीही $Z_1$, $Z_2$, \dots{}
  अशी यादी दिली असता, यादीत नसेल असा संख्यांचा संच~$\overline{Z}$
  रचता येतो, हे दाखवले. तो यादीत नसतो याची खात्री यामुळे मिळते की
  प्रत्येक $n \in
  \PosInt$ साठी $n \in Z_n$ तेव्हा आणि केवळ तेव्हाच
  $n \notin \overline{Z}$. त्यामुळे $Z_n$ आणि $\overline{Z}$ यांपैकी
  एकाचा !!a{element} असलेली, पण दुसऱ्याचा नसलेली कोणती तरी संख्या
  नेहमी असते. येथेही तीच कल्पना वापरतो; फरक एवढाच की निर्देशांक~$n$
  आता~$\PosInt$ चे नसून~$A$ चे !!{element}s आहेत. $\overline{A}$ ची
  व्याख्या अशी केली आहे की प्रत्येक $x \in A$ साठी तो~$g(x)$ पेक्षा
  वेगळा असतो, कारण $x \in g(x)$ तेव्हा आणि केवळ तेव्हाच
  $x \notin \overline{A}$. पुन्हा, $g(x)$ आणि $\overline{A}$ यांपैकी
  एकाचा !!a{element} असलेला, पण दुसऱ्याचा नसलेला~$A$ चा कोणता तरी
  !!a{element} नेहमी असतो. म्हणूनच जसा $\overline{Z}$ हा $Z_1$,
  $Z_2$, \dots{} या यादीत येऊ शकत नाही, तसा $\overline{A}$ हा~$g$ च्या
  व्याप्तीत येऊ शकत नाही.}{}

\oliflabeldef{sfr:siz:nen-alt:thm:nonenum-pownat}{
  \olref{thm:cantor} ची सिद्धता आणि
  \olref[nen-alt]{thm:nonenum-pownat} ची सिद्धता यांची तुलना करणे
  उद्बोधक आहे. तेथे~$\Nat$ च्या उपसंचांची कोणतीही $N_0$, $N_1$,
  $N_2$, \dots{} अशी यादी दिली असता, यादीत नसेल असा संख्यांचा संच~$D$
  रचता येतो, हे दाखवले. प्रत्येक $n \in \Nat$ साठी $n \in N_n$ तेव्हा
  आणि केवळ तेव्हाच $n \notin
  D$ असल्यामुळे तो यादीत नसतो याची खात्री मिळते. येथेही तीच कल्पना
  वापरतो; फरक एवढाच की निर्देशांक~$n$ आता~$\Nat$ चे नसून~$A$ चे
  !!{element}s आहेत. $D$ ची व्याख्या अशी केली आहे की प्रत्येक $x
  \in A$ साठी तो~$g(x)$ पेक्षा वेगळा असतो, कारण $x \in g(x)$ तेव्हा
  आणि केवळ तेव्हाच $x \notin D$.}{}

या सिद्धतेची रसेल यांच्या विरोधापत्तीच्या सिद्धतेशीही तुलना करणे उपयुक्त
आहे: \olref[sfr][set][rus]{thm:russells-paradox}. खरे तर, कँटर यांच्या
प्रमेयातूनच रसेल यांना स्वतःची विरोधापत्ती सुचली.
\end{explain}

\begin{prob}
  कोणत्याही संच~$A$ साठी $g\colon \Pow{A} \to
  A$ असे !!a{injection} असू शकत नाही, हे दाखवा. सूचना: $g\colon
  \Pow{A} \to A$ हे !!{injective} आहे, असे समजा. पुढील संच विचारात घ्या:
  $D = \Setabs{g(B)}{B \subseteq A \text{ आणि
  } g(B) \notin B}$. $x = g(D)$ असे घ्या. $g$ हे !!{injective} आहे हा
  हे तथ्य वापरून व्याघात मिळवा.
\end{prob}

\end{document}
